% =====================================================================
%  Thème 3 — EXERCICES — Niveau FACILE
% =====================================================================
\documentclass[11pt,a4paper]{article}
\usepackage[utf8]{inputenc}
\usepackage[T1]{fontenc}
\usepackage{lmodern}
\usepackage[french]{babel}
\usepackage{amsmath,amssymb,mathtools}
\usepackage{icomma}
\usepackage{xcolor}
\usepackage[most]{tcolorbox}
\usepackage{enumitem}
\usepackage{array}
\usepackage{booktabs}
\usepackage{fancyhdr}
\usepackage[hidelinks]{hyperref}
\usepackage[a4paper,margin=2.2cm,headheight=26pt,headsep=10pt]{geometry}

\definecolor{primary}{RGB}{223,87,99}
\definecolor{primarydark}{RGB}{150,42,55}

\tcbset{boxbase/.style={enhanced,breakable,boxrule=0.9pt,arc=2.5pt,
   left=3mm,right=3mm,top=2.3mm,bottom=2.3mm,
   fonttitle=\bfseries,coltitle=white,toptitle=0.6mm,bottomtitle=0.6mm}}
\newtcolorbox[auto counter]{exercice}[1][]{boxbase,colback=primary!4,colframe=primary,
  title={\textsc{Exercice}~\thetcbcounter\ifx\relax#1\relax\relax\else\textmd{ --- \itshape #1}\fi}}

\pagestyle{fancy}
\fancyhf{}
\fancyhead[L]{\small\textcolor{primarydark}{\textbf{Fiche 6} — Les limites}}
\fancyhead[R]{\small\textcolor{primarydark}{\textsc{Thème 3 — Facile}}}
\fancyfoot[C]{\small\thepage}
\renewcommand{\headrulewidth}{0.8pt}
\renewcommand{\headrule}{\hbox to\headwidth{\color{primary}\leaders\hrule height \headrulewidth\hfill}}

\newcommand{\R}{\mathbb{R}}

\begin{document}

\begin{center}
{\Large\bfseries\color{primarydark}Thème 3 — Limites à l'infini : le terme dominant}\\[2pt]
{\large\color{primary}Exercices — Niveau Facile}
\end{center}
\vspace{4pt}
{\small\itshape But : ne garder que le terme de plus haut degré, et lire le résultat (signe, parité,
comparaison de degrés).}
\vspace{6pt}

\begin{exercice}[monômes à l'infini]
Donner la limite ($+\infty$ ou $-\infty$).
\begin{enumerate}[label=\alph*),leftmargin=2em,itemsep=3pt]
  \item $\displaystyle\lim_{x\to+\infty} x^3$ \quad b) $\displaystyle\lim_{x\to-\infty} x^3$
  \item $\displaystyle\lim_{x\to-\infty} x^4$ \quad d) $\displaystyle\lim_{x\to-\infty} (-x^2)$
\end{enumerate}
\end{exercice}

\begin{exercice}[polynômes : garder le terme dominant]
Calculer en ne conservant que le terme de plus haut degré.
\begin{enumerate}[label=\alph*),leftmargin=2em,itemsep=3pt]
  \item $\displaystyle\lim_{x\to+\infty}(2x^3-5x^2+7)$
  \item $\displaystyle\lim_{x\to-\infty}(x^2+10x+1)$
  \item $\displaystyle\lim_{x\to+\infty}(-4x^5+x)$
\end{enumerate}
\end{exercice}

\begin{exercice}[rationnelles : degrés égaux]
Chaque quotient a un numérateur et un dénominateur de \emph{même} degré. Donner la limite.
\begin{enumerate}[label=\alph*),leftmargin=2em,itemsep=3pt]
  \item $\displaystyle\lim_{x\to+\infty}\frac{3x^2+1}{5x^2-x}$
  \item $\displaystyle\lim_{x\to-\infty}\frac{-2x+7}{4x+3}$
\end{enumerate}
\end{exercice}

\begin{exercice}[rationnelles : degré du bas plus grand]
Ici le dénominateur « gagne » : la limite est nulle. Vérifier.
\begin{enumerate}[label=\alph*),leftmargin=2em,itemsep=3pt]
  \item $\displaystyle\lim_{x\to+\infty}\frac{2x+1}{x^2+3}$
  \item $\displaystyle\lim_{x\to-\infty}\frac{5}{x^3}$
\end{enumerate}
\end{exercice}

\begin{exercice}[rationnelles : degré du haut plus grand]
Ici le numérateur « gagne » : la limite est infinie. Donner le signe.
\begin{enumerate}[label=\alph*),leftmargin=2em,itemsep=3pt]
  \item $\displaystyle\lim_{x\to+\infty}\frac{x^3+1}{2x+1}$
  \item $\displaystyle\lim_{x\to+\infty}\frac{-x^4}{3x^2+1}$
\end{enumerate}
\end{exercice}

\end{document}
